
%=======================================================================================
\chapter{Data}
\label{ch:data}

\sr{This chapter describes the data used for model training and evaluation, including the study area, cohort construction and retained variables. Table~\ref{tab:data_summary} summarises the predictor and target variables, their sources and temporal coverage. Feature engineering, data splits and modelling decisions are described in Chapter~\ref{ch:methodology}.}


%-
\section{Study area and cohort coverage}
\label{sec:data-source}

\sr{The dataset covers approximately $16\text{~km} \times 12\text{~km}$ of Aberfoyle and contains surveys from \{2002, 2006, 2008, 2012, 2021, 2023\}. Each of the 795,645 raw rows represents one plot in one survey year and contains airborne laser scanning (ALS)-derived height and canopy variables alongside stand and management fields (Table~\ref{tab:data-fields}).} Most plots are a $20\text{~m} \times 20\text{~m}$ square; however, \sr{28.8\% are clipped at compartment or sub-compartment boundaries (Figure~\ref{fig:data-study-area}). The analysis was restricted to evergreen Sitka spruce (SS), the most common and widely distributed species in Aberfoyle, to reduce seasonal canopy variation between surveys.} \lim{Discuss unknown flight dates, sensor differences and remaining windthrow or abiotic disturbance in the final limitations section.}

\begin{figure}[!htbp]
\centering
\includegraphics[width=0.95\textwidth]{figures/fig_data/aberfoyle_area.png}
\caption{Aberfoyle study area, spatial hierarchy, and plot geometry.}
\label{fig:data-study-area}
\end{figure}

\app{Full raw column inventory (all 38 fields in \texttt{LiDAR\_Years}), including fields never
used downstream (\texttt{flyr}, constant at zero for every row; \texttt{block}, a
99.999\%-identical duplicate of \texttt{blk}).}


%-
\section{Cohort construction and coverage}
\label{sec:data-cohorts}

Due to a jump in LiDAR scan coverage between 2006 and 2008 (an increase of 37~$\text{km}^2$), two cohorts were created: a four-survey cohort ($\in \{2008, 2012, 2021, 2023\}$) and a six-survey cohort ($\in \{2002, 2006, 2008, 2012, 2021, 2023\}$). 

\sr{The six-survey cohort is nested within the four-survey cohort (Figure~\ref{fig:cohort_coverage_funnel}). The four-survey cohort provides broader spatial coverage, whereas the six-survey cohort provides longer trajectories for temporal evaluation and sensitivity analysis.} \fc{Explain why each cohort is assigned to the spatial, temporal and sensitivity evaluations in Methodology.}

\sr{Sequential filters required complete survey coverage, pure Sitka spruce stands, a planting year after 1850 and before the survey, stand age between $0$ and $200$ years, height between $0\text{~m}$ and $60\text{~m}$, and stand age of at least 30 years in 2023 (Table~\ref{tab:data-cleaning-funnel}). The height bound follows the local ecological maximum for Sitka spruce in Aberfoyle \cite{mason2011Sitka_C1_SS}; the maturity filter excludes young stands with less reliable top-height measurements \cite{suarezminguez2008Practical_C1_SS}. The planting-year check removed 629 four-survey plots (0.87\% of Sitka spruce plots), while the age and height checks removed no further plots. Most attrition came from requiring complete survey coverage and restricting the data to Sitka spruce.}


\begin{table}[!htbp]
\centering
\scriptsize
\begin{tabular}{lrr}
\toprule
Stage & 4-survey & 6-survey \\
\midrule
0. Raw dataset               & 230,359 & 230,359 \\
1. Cohort years present      & 230,345 & 230,359 \\
2. Complete survey coverage  & 112,782 &  18,781 \\
3. Sitka spruce only         &  72,395 &  13,897 \\
4. Valid planting year       &  71,766 &  13,897 \\
5. Plausible age (0-200y)   &  71,766 &  13,897 \\
6. Plausible height (0-60m) &  71,766 &  13,897 \\
7. Mature stands only (Age $\ge 30$ in 2023) & \textbf{58,112} & \textbf{13,769} \\
\bottomrule
\end{tabular}
\caption{\sr{Unique plots remaining after each data filter. Geographic coverage is shown in Figure~\ref{fig:data-study-area}.}}
\label{tab:data-cleaning-funnel}
\end{table}

%-
\section{Target variable and original stand fields}
\label{sec:data-response}

\sr{The response was the 95th percentile ALS return height above ground (\texttt{elev\_percentile\_95th}), averaged within each plot boundary for each survey.} \fc{Keep the justification for using this measure as top height, and its relation to yield class and volume, in Background.} The survey records also include planting year, management history \sr{and Forest Research yield and volume estimates}. Table~\ref{tab:data-fields} lists the different variables and their role in the study. \sr{Fields derived from the height response were excluded from prediction to prevent target leakage.} \fc{Confirm whether \texttt{yldc} is derived from this response; Table~\ref{tab:data-fields} currently says that it is not.}

\begin{table}[!htbp]
\scriptsize
\setlength{\tabcolsep}{3pt}
\renewcommand{\arraystretch}{0.95}
\begin{tabular}{p{2.6cm}p{7.0cm}p{3.2cm}}
\hline
Field & Meaning or derivation & Role in this study \\
\hline
\texttt{elev\_percentile\_95th} & 95th percentile of ALS return heights above ground, used without correction & Response variable \\
\texttt{Top\_Height95} & $=\texttt{elev\_percentile\_95th}\times1.1$; Forest Research's own corrected height estimate & Retained for reporting but not modelled \\
\texttt{plyr}, \texttt{Age} & Planting year; \texttt{Age} = survey year $-$ \texttt{plyr} & Input feature; used for data cleaning \\
\texttt{yldc} & An existing Forest Research yield-class field, not calculated from this project's LiDAR-derived height response & Excluded from prediction; used as a comparison variable \\
\texttt{whcl} & Ordinal Forest Research windthrow hazard class, 0-6 & Management descriptor \\
\texttt{CanopyCover} & Canopy-cover fraction from the same ALS survey and processing step as the height response & Input feature \\
\texttt{Thin}, \texttt{last\_thinn}, \texttt{next\_thin\_} & Thinning history; \texttt{Thin} is binary (1 if the stand has been thinned, 0 if never thinned); \texttt{last\_thinn}/\texttt{next\_thin\_} record the last and next thinning years & Input feature \\
\texttt{Vol95}, \texttt{area}, \texttt{g1}, \texttt{g3} & Forest Research volume estimate, calculated from height, plot area and canopy cover & Used for disturbance screening \\
\texttt{GYCspec95}, \texttt{p1-p5} & General Yield Class, calculated directly from a plot's own measured height (\texttt{Top\_Height95}) \ct{Forest Research, "Use of Airborne Laser Scanning (ALS) for Forest Inventory" - source of p1-p5, full reference not yet confirmed} & Excluded from prediction because its calculation
incorporates the height response, creating target leakage \\
\hline
\end{tabular}
\caption{\sr{Target variable and original stand and management fields retained in the master export.}}
\label{tab:data-fields}
\end{table}
%-
\section{Derived terrain, wind, climate and soil variables}
\label{sec:data-environmental}

Each plot was matched with external terrain, wind, climate, and soil variables (Table~\ref{tab:data-env-sources}). \sr{Full variable definitions are provided in Appendix~\ref{tab:data-env-variables_detail_app}.}

\begin{table}[!htbp]
\tiny
\setlength{\tabcolsep}{2pt}
\renewcommand{\arraystretch}{0.9}
\begin{tabular}{p{1.6cm}p{2.1cm}p{1.6cm}p{1.5cm}p{2.3cm}p{1.6cm}p{2.3cm}p{0.7cm}}
\hline
Data family & Source & Resolution & Plots/cell (obs.; geom.\ max) & Temporal coverage & Static/ time-varying & Main purpose & Citation \\
\hline
Forest plots & Forest Research & $\sim$20$\times$20\,m & Native unit & 2002-2023 (6 years) & Time-varying & Response, stand fields & - \\
Terrain & OS Terrain 50, Ed.\ 10 & 50\,m grid & 4.3; 6.25 & Ed.\ 10; flown 2023-05-29, created 2025-12-25 (tile NN40 only; 5 other tiles unconfirmed) & Static & Elevation, slope, aspect & \cite{ordnancesurveyTerrain50_C3_LID} \\
Soil (categorical) & CEH nat.\ capital rasters & 50\,m grid & 4.3; 6.25 & 2020 (confirmed) & Static & Pedotope$^{\dagger}$, drainage, texture & \cite{feeney2026CEHNaturalCapital_C3_LID} \\
Soil pH & SoilGrids v2.0, ISRIC & 250\,m grid & 53; 156.25 & Released May 2020 & Static (0-5\,cm) & Soil acidity & \cite{poggio2021SoilGrids_C3_LID} \\
Wind climatology & Global Wind Atlas & 250\,m grid (bulk); $\sim$155$\times$278\,m (10\,m-height REST variant) & n/a (anisotropic for REST variant) & 10-yr climatology, 2006-2015 & Static, long-term & General wind exposure & \cite{davis2023GlobalWindAtlas_C3_LID} \\
Observed wind & MIDAS Open, 3 stns & Point & Distance: 30.7/47.7/50.5\,km & 2002-2023, per interval & Time-varying & Storm/wind conditions & \cite{metoffice2019MIDASOpen_C3_LID} \\
Climate (long-term) & CHELSA v2.1 BIOCLIM+ & 1\,km grid & 482; 2{,}500 & 1981-2010 & Static & GDD, precipitation & \cite{karger2017CHELSA_C3_LID} \\
Climate (matched) & HadUK-Grid & 1\,km grid & 482; 2{,}500 & Cohort-year-averaged & Const.\ within trajectory & Temperature, frost & \cite{hollis2019HadUKGrid_C3_LID} \\
Climate (excluded) & ERA5-Land & 0.1$^\circ$ ($\sim$9\,km native) & - & Reanalysis, ongoing & Static/reanalysis & Excluded: 8 distinct values only & \cite{munozsabater2021ERA5Land_C3_LID} \\
Spatial position & OS Open Roads/Rivers & Vector & Native & Refreshed 6-monthly; download date unlogged & Static & Edge/context effects & \cite{ordnancesurveyOpenRoads_C3_LID, ordnancesurveyOpenRivers_C3_LID} \\
\hline
\end{tabular}
\caption{\sr{External data sources matched to each plot by point extraction at its coordinate centre. An estimated centre was used for clipped plots.} \fc{Describe the extraction procedure in Methodology and check whether ``MAUP technique'' is the correct term; MAUP describes a source of spatial bias rather than a method for calculating a centre.}}
\label{tab:data-env-sources}
\end{table}

\sr{The three wind sources represent different aspects of exposure: long-term modelled wind climate (Global Wind Atlas), terrain shelter (TOPEX), and observed conditions at three MIDAS stations.} \fc{Keep the forestry justification for these sources in Background and the interval summarisation procedure in Methodology.} \lim{Discuss the distance between Aberfoyle and the MIDAS stations in the final limitations section.}

\sr{CHELSA v2.1 is a downscaled global climate dataset whose BIOCLIM+ variables represent long-term 1981--2010 conditions \cite{karger2017CHELSA_C3_LID}. HadUK-Grid is a gridded UK observational dataset; its temperature and ground-frost variables were averaged over the years represented by each cohort \cite{hollis2019HadUKGrid_C3_LID}.} \fc{Explain the temporal aggregation procedure and why variables were not matched to individual survey years in Methodology.} \lim{Discuss the inability of static or cohort-averaged variables to represent individual growth intervals in the final limitations section.}

\fc{Move the aspect transformation to Methodology: terrain aspect was converted to northness ($\cos\theta$) and eastness ($\sin\theta$) to avoid the discontinuity between $0^\circ$ and $360^\circ$ \cite{moisen2002Comparing_C3_LID}.} \app{Full derivation and frost-hollow threshold.}


%-
\lim{Move the former ``Strengths and limitations of the data'' subsection to the final limitations section after Results and Discussion. Cover: the single-forest scope; observational and non-causal attribution; incomplete management records; inherited bias and smoothing in modelled environmental products; the coverage--trajectory trade-off between cohorts; and plot-level rather than individual-tree inference. Keep the Conclusion brief and derive external-site validation and finer local measurements as future work.}
